

\section{Common Authentication Vulnerabilities}\label{sec:common-vulnerabilities}
Bad actors find the most creative ways to attack authentication systems, these are called attack vectors. The goal of this section is to outline how the vulnerabilities work, which the paper will address in later sections, and also to shine a ligh on why traditional password or even basic MFA schemes struggle against them.s


\subsection{Weak Passwords and Brute-Force Attacks}
Users often choose simple passwords, making them easy targets for brute-force attacks \cite{Ma2022ProspectsPassword}. A simple brute-force attack tries to crack the user's password by trying random character combinations. Rate limiting and sufficient password length makes the attack redundant. A dictionary attack is a more sophisticated version of a brute-force attack, where attackers use common wordlists, lists of frequenly used passwords or they systematically guess combinations to find the password. In the case of brute-force not cryptographic weaknesses, but bad password hygiene and the site's lack of password complexity enforcement leads to the compromise,.


\subsection{Credential Stuffing and Password Reuse}
Credential stuffing exploits password reuse across websites. In the event of a databreach the user's password may get leaked, this creates an opportunity for bad actors, who then use the credentials on other sites to try to gain access to other accounts.\cite{Nisenoff2023PasswordReuse}. The process can be well automated, thus the reuse of credentials turns one breach into many.


\subsection{Social Engineering Attacks}
Many attacks rely on human gullability instead of technical flaws. Phishing, SIM swap, and push fatigue attacks are such attacks.


\subsubsection{Phishing and OTP Relay}
Phishing deceives users into revealing credentials on fake websites, that mimic the working of the site the user is familiar with. Attackers can proxy the login and capture OTP codes in real time, bypassing traditional 2FA, this is commonly referred to as a real-time phishing proxy. OTP capturing is Adversary-in-the-Middle type attack\cite{Kasemodel2025BlueTOTP}.

\subsubsection{SIM Swap Attacks}
The way SMS based MFA works is that the user gets an SMS containing a secret code, generated by the server, this while may seems bulletproof requires complete trust in the user's telecommunication service provider. The attackers may be able to impersonate victims and hijack their phone number (with the service provider unknowingly complicit), thus gaining control of SMS-based OTPs\cite{Lee2020SIMswap}.


\subsubsection{MFA Push Fatigue}
Push-based MFA can be exploited by flooding users with repeated prompts until they approve one, often by mistake or frustration \cite{CISA2023PhishingResistantMFA}.

\subsection{Server-Side Data Breaches}
When there is a data breach in the authentication server's database, password hashes or shared secrets (like TOTP seeds) can get exposed, and then the attackers can crack these sensitive information offline to gain access to the accounts \cite{OWASPPasswordStorage}.

\subsection{Adversary-in-the-Middle tampering}
Interception and modification of traffic between the client and server by a third party falls into the category of tampering. One would say, that TLS provides transport security, but by itself it does not prevent manipulation if it gets compromised or stripped \cite{OWASPMitM}.

\subsection{Denial of Service via Account Lockout}
Sometimes security measures themselves can be weaponized against the users or the service, attackers trigger account lockouts by repeatedly submitting invalid credentials for a target user. This exploits the lockout policy, while it keeps the user safe from brute-force attacks, it can lead to them not being able to access their account \cite{NISTSP80063B}.